\section{Numerical scheme}
\label{sec:numerics}

\subsection{The radial grid}
\label{sec:numerics:grid}

The domain above the base level $r = 1$ is discretized into $N$ finite-volume cells
(\texttt{Grid cells:}, default 500) with $N_g = 2$ ghost cells at each end;
all grid arrays run over $1-N_g \le j \le N+N_g$.  \texttt{define\_grid.f90}
builds the cell centers $r_j$, the interfaces $r_{j+1/2}$ as their midpoints,
and the widths $\Delta r_j = r_{j+1/2} - r_{j-1/2}$.  The last is an exact
identity of the discretization and not an approximation: it is the same face
pair that forms $\Delta V_j$ in \eqref{eq:geom}, the same pair the
gravitational source differences, and the same width the column integrals
$n_j\,\Delta r_j$ and the cell optical depths use.

Three grid types are available, and since 2026-09-27 every one of them
places the first face exactly on the base level, $r_{\rm edg}(0) = 1$, so
that the reservoir stated there is the state of that face
(Sec.~\ref{sec:bcic:base}); the two lower ghost cells lie below the level.
\texttt{Uniform} takes $\Delta r = (r_{\max}-1)/N$ and
$r_j = 1 + (j - \tfrac12)\Delta r$, so its faces are $1 + j\,\Delta r$ and the
outer face of cell $N$ is $r_{\max}$.  \texttt{Stretched} places the centers
in geometric progression, $r_j = 2x^j/(1+x)$, so that the face between
$r_0$ and $r_1$ is the level, with $x$ the root of
$(N+N_g)\ln x - \ln[(1+x)/2] - \ln r_{\max} = 0$ (the center of the outermost
ghost at $r_{\max}$), found by Newton's method from
$x = r_{\max}^{1/(N+N_g)}$, which climbs to it monotonically because the
function is increasing and concave.  \texttt{Mixed}, the type every
shipped configuration uses, lays $N_{\rm low}$ uniform cells of width
$\Delta r_{\rm base}$ on the base and then stretches the remaining
$N_{\rm up} = N - N_{\rm low}$ cells by a factor $k$ that solves
%
\begin{equation}
  \frac{1 - k^{N_{\rm up}}}{1-k}
  = \frac{r_{\max} - r_{N_{\rm low}}}{\Delta r_{\rm base}}
  \label{eq:stretch}
\end{equation}
%
by Newton-Raphson to $10^{-8}$, which is Eq.~(10) of \citet{Caldiroli2021}.
The two numbers are the key \texttt{Base grid [dr,cells]:} and default to
$\Delta r_{\rm base} = 2\times10^{-4}\,R_p$ and $N_{\rm low} = 50$, the values
ATES hard-codes.  The widths are then smoothed once by a 1-2-1 pass running
downward from the junction to the base, the centers are rebuilt from the
smoothed widths, and the grid is mapped by the affine map that sends its
first face to $1$ and the center of its outermost ghost to $r_{\max}$, which
multiplies every width by one factor and so keeps the smoothed shape.  For
all three types the first ghost center is then restated as the mirror of the
first cell center, $r_0 = 2 - r_1$, which moves it by rounding only and makes
the first face exactly $1.0$; the run stops if that face is not 1 or if a
ghost cell reaches $r \le 0$.  Until 2026-09-27 the level was a ghost
\emph{center} ($r_0 = 1$ on the \texttt{Mixed} grid, $r_{1-N_g} = 1$ on the
other two), so every cell center has since moved (by half a cell on the
\texttt{Mixed} grid, by one and a half on the other two) and a state
written before then does not load on the present grid
(Sec.~\ref{sec:bcic:restart}).
The requirement behind $\Delta r_{\rm base}$ is that the pressure scale
height $H = k_BT/(\bar m g)$, with $\bar m = \rho/n_{\rm tot}$ the mean
particle mass, be resolved by several cells at the base; a
configuration that violates it is the one that produces a base-layer momentum
imbalance no solver removes.

Two windows are located once, at grid construction: $j_{\min}$, the first
cell with $r \ge r_{\rm esc}$ (\texttt{Escape radius [R\_p]:}), which the
marching convergence measure of Sec.~\ref{sec:numerics:du} is taken over, and
$j_{\rm flux}$, the first cell with $r \ge r_{\rm flux}$ (default
$1.2\,R_p$), which the stationary flux gate uses.  Both are clamped to
mid-domain, with a message, where the requested radius lies outside the
domain: a deep Roche-lobe-filling case can have $r_{\max}$ smaller than the
default escape radius, and an empty window would make the measure $-\infty$
and the state converged at step zero.

\subsection{Reconstruction}
\label{sec:numerics:reconstruction}

Both schemes reconstruct the \emph{primitive} variables $W = (\rho, v, p)$ to
the two sides of every face, and both write the left state of face $j$ from
cell $j$ and the right state of face $j$ from cell $j+1$.

\paragraph{PLM.}
\texttt{PLM\_rec.f90} builds a limited slope from the three cell averages
$j-1, j, j+1$ and evaluates it at the two faces,
%
\begin{equation}
  W_{L,j+1/2} = W_j + \tfrac12\sigma\,(r_{j+1}-r_j),
  \qquad
  W_{R,j-1/2} = W_j - \tfrac12\sigma\,(r_j-r_{j-1}),
  \label{eq:plm}
\end{equation}
%
with $\sigma$ the generalized minmod (monotonized central) slope of the right,
left and central difference quotients at $\theta = 2$,
%
\begin{equation}
  \sigma = \mathrm{minmod}\!\left(
     \theta\,\frac{\Delta W_{j+1/2}}{r_{j+1}-r_j},\;
     \theta\,\frac{\Delta W_{j-1/2}}{r_j-r_{j-1}},\;
     \frac{W_{j+1}-W_{j-1}}{r_{j+1}-r_{j-1}}\right),
  \label{eq:minmod}
\end{equation}
%
the slope being zero whenever the three arguments do not share a sign.  This
is the generalized minmod limiter as \citet{Kaeppeli2016} state it, at the
least diffusive value for which the scheme stays total-variation diminishing;
ATES uses the same limiter and attributes it to \citet{Kurganov2000}.  The
difference from ATES is that the quotients here are formed on the actual
center spacing of a stretched grid.  The same scalar limiter is exported as
\texttt{PLM\_rec\_scalar}, so that a transported mass fraction rides on the
same discretization as the density whose face flux carries it.

\paragraph{WENO3.}
\texttt{Reconstruction.f90} uses the third-order WENO reconstruction with the
nonlinear weights of \citet{Yamaleev2009}.  With
$\Delta W_{j\pm1/2}$ the two interface jumps bounding cell $j$,
%
\begin{equation}
  \tau = \Delta W_{j+1/2} - \Delta W_{j-1/2},\qquad
  S_0 = 1 + \frac{\tau^2}{\Delta W_{j+1/2}^2 + \Delta r_j^2},\qquad
  S_1 = 1 + \frac{\tau^2}{\Delta W_{j-1/2}^2 + \Delta r_j^2},
  \label{eq:weno_weights}
\end{equation}
%
which is their $1 + \tau/(\epsilon + \beta_r)$ with their already-squared
$\tau$ (their Eq.~22) and the floor $\epsilon = \Delta r_j^2$ folded in.  The
face states are
%
\begin{equation}
  \begin{aligned}
    W_{L,j+1/2} &= W_j + \frac{S_0\,C_2(j)\,\Delta W_{j+1/2}
                    + D_1(j)\,S_1\,C_1(j-1)\,\Delta W_{j-1/2}}
                   {S_0 + D_1(j)\,S_1},\\
    W_{R,j-1/2} &= W_j - \frac{D_2(j)\,S_0\,C_2(j)\,\Delta W_{j+1/2}
                    + S_1\,C_1(j-1)\,\Delta W_{j-1/2}}
                   {D_2(j)\,S_0 + S_1},
  \end{aligned}
  \label{eq:weno_faces}
\end{equation}
%
with the purely geometric coefficients
$C_1(j) = \Delta V_{j+1}/(\Delta V_j + \Delta V_{j+1})$, $C_2 = 1 - C_1$,
$D_1(j) = \Delta V_{j+1}/(\Delta V_j + \Delta V_{j-1})$ and
$D_2(j) = \Delta V_{j-1}/(\Delta V_j + \Delta V_{j+1})$.  Each candidate is
the linear interpolant of two cell averages \emph{in the volume coordinate},
so the jump to a neighbor is weighted by the cell's own volume share; these
shares, and not the nonlinear weights, are what set the order on a stretched
grid.  The ideal weights $d_0 = 2/3$, $d_1 = 1/3$ of the uniform-grid scheme
enter only through their ratio, which is what $D_1$ and $D_2$ carry: exactly
$1/2$ on a grid uniform in the volume coordinate and $1/2 + O(h)$ otherwise.

Two limits of the scheme are stated at the routine.  What runs is a
third-order WENO reconstruction \emph{with} those weights and not the
energy-stable ESWENO scheme of the same paper, whose extra dissipation flux
(their Eq.~48) is not computed here.  And the floor $\epsilon = \Delta r_j^2$
has the required $h^2$ scaling but drops the solution scale of the published
prescription, so it is not invariant under a rescaling of the solution; what
that costs is the stencil biasing in the tenuous outer domain and not the
order, which is flat over four decades of the coefficient.

\paragraph{Positivity of the face states.}
Neither scheme has a positivity property: a face value of $\rho$ or $p$ can
cross zero while every cell average on the stencil is positive, and at Mach
60 the thermal pressure is $2\times10^{-4}$ of the total energy density, so a
reconstruction across a jump tips it negative and the sound speed takes the
square root of it.  \texttt{positivity\_limited\_faces} therefore scales each
offending face state \emph{continuously} toward its own cell average,
%
\begin{equation}
  W_{\rm face} \leftarrow W_{\rm avg} + \theta\,(W_{\rm rec} - W_{\rm avg}),
  \qquad
  \theta = \min_{q \in \{\rho,\,p\}}
           \frac{q_{\rm avg} - \varepsilon}{q_{\rm avg} - q_{\rm rec}},
  \qquad \varepsilon = \epsilon_{\rm mach}\,q_{\rm avg},
  \label{eq:posface}
\end{equation}
%
which is the linear scaling idea of \citet{Zhang2010} applied to the
primitive vector, and not their limiter: they scale the conserved vector in
two steps with $\theta$ the root of a quadratic, one $\theta$ per cell against
one per face state here, and their floor is absolute against one unit in the
last place of the cell average here.  What carries over is their Lemma 2.5,
that a convex combination of an admissible average with the reconstruction is
admissible; their Theorem 2.1 does not, its hypotheses being a Gauss-Lobatto
quadrature set and a CFL condition nothing here imposes.  Continuity is the
point: a hard switch at the crossing is a step discontinuity of the residual
and no stationary solve can descend across one.  A face whose reconstruction
is admissible has $\theta = 1$ and is not rewritten.

\paragraph{The PLM to WENO3 continuation.}
The two-stage recipe \texttt{Reconstruction scheme: PLM+WENO3} marches with
PLM until the convergence measure falls below the first \texttt{du\_th} value
(or its plateau is detected) and then continues with WENO3.  Replacing one
discretization by another between two time steps is a jump of the right-hand
side, so the key \texttt{Reconstruction continuation:} interpolates instead.
\texttt{reconstruction\_continuation\_rhs} (\texttt{steady\_residual.f90})
evaluates the homotopy
%
\begin{equation}
  R_\lambda(u) = (1-\lambda)\,R_{\rm PLM}(u) + \lambda\,R_{\rm WENO3}(u),
  \label{eq:homotopy}
\end{equation}
%
returning its flux-difference and source parts separately, and blends the
stored interface fluxes, the face pressures and the three momentum terms on
the same $\lambda$.  Everything the two schemes disagree about is inside the
two evaluations: the face reconstruction, the extrapolated reconstructed-face
boundary states, and the pressure split of Sec.~\ref{sec:hydro:pressure}.
At $\lambda \le 0$ and $\lambda \ge 1$ a single evaluation is taken and the
two pure operators are reproduced exactly; only the open interval costs two
right-hand sides.  Because the admissible set $\rho>0$, $\rho e>0$ is convex
and each stage is then a convex combination of two forward Euler stages, the
continuation cannot manufacture a positivity failure that neither scheme has.
The ramp is adaptive: $\lambda$ advances by $\Delta\lambda$ where the relative
state change of the step stays below a tolerance and $\Delta\lambda$ is halved
where it does not.

\subsection{Approximate Riemann solvers}
\label{sec:numerics:riemann}

\texttt{Num\_Fluxes.f90} carries three fluxes, selected by the mandatory key
\texttt{Numerical flux:}.  Each side's total energy and sound speed are built
with its \emph{own} heat capacity, which is what the caloric equation of state
makes different across a face inside the $\mathrm{H}_2$ front; the physical
flux is
%
\begin{equation}
  F(W) = \bigl(\rho v,\ \rho v^2 [+\,p],\ v(\mathcal{E}+p)\bigr),
  \label{eq:physflux}
\end{equation}
%
the bracketed pressure being present under PLM without the well-balanced
option, as Sec.~\ref{sec:hydro:pressure} requires.

\paragraph{HLLC.}
The three-wave solver of \citet{Toro1994}, with the Davis/Einfeldt speed
bounds of \citet{Batten1997},
%
\begin{equation}
  S_L = \min\bigl(0,\, v_L - a_L,\, v_R - a_R\bigr),\qquad
  S_R = \max\bigl(0,\, v_L + a_L,\, v_R + a_R\bigr),
  \label{eq:hllc_speeds}
\end{equation}
%
\begin{equation}
  S_* = \frac{p_R - p_L + \varphi_L v_L - \varphi_R v_R}{\varphi_L - \varphi_R},
  \qquad \varphi_K = \rho_K (S_K - v_K),
  \label{eq:sstar}
\end{equation}
%
and the star states
%
\begin{equation}
  u^*_K = \frac{\rho_K(S_K - v_K)}{S_K - S_*}
  \left(1,\ S_*,\
        \frac{\mathcal{E}_K}{\rho_K} + (S_* - v_K)\Bigl[S_* + \frac{p_K}{\varphi_K}\Bigr]
  \right)^{\!\top},
  \label{eq:hllc_star}
\end{equation}
%
the flux being $F_K$, $F_K + S_K(u^*_K - u_K)$ or the mirror of those
according to the signs of $S_L$, $S_*$, $S_R$.  Because $S_L$ and $S_R$ are
clamped at zero, a face with every physical wave moving outward returns
$F(W_L)$ and never forms an expression in the right state.  The face pressure
$p_{\rm out}$ that the momentum assembly of \eqref{eq:dF2} needs is the
pressure of whichever side the branch took.  HLLC is the flux every shipped configuration under \texttt{examples/} and
\texttt{backup/regression/} states; the key is mandatory and has no default.
It is not the flux of every run: the seven lowest-XUV cases of the LHS 1140 b
catalog state \texttt{ROE}, because on a 500-cell grid whose base is subsonic
to $M\sim10^{-7}$ the HLLC flux does not reach the stationary root there. That
is a resolution statement about the flux and not about the wind: the same
unmodified HLLC flux certifies the same case on 1000 cells, and the state it
reaches is the 500-cell Roe state to 0.1--3 per cent above $1.005\,R_p$
(measured 2026-09-17).

\paragraph{Roe.}
The linearized solver of \citet{Roe1981}, with the averages
%
\begin{equation}
  \tilde\rho = \sqrt{\rho_L\rho_R},\qquad
  \tilde v = s v_L + (1-s) v_R,\qquad
  \tilde H = s H_L + (1-s) H_R,\qquad
  s = \frac{\sqrt{\rho_L}}{\sqrt{\rho_L}+\sqrt{\rho_R}},
  \label{eq:roe_avg}
\end{equation}
%
$\tilde a = \sqrt{(\bar\gamma-1)(\tilde H - \tilde v^2/2)}$, and the flux
%
\begin{equation}
  F_{j+1/2} = \tfrac12\bigl(F_L + F_R\bigr)
            - \tfrac12\sum_{i=1}^{3}\tilde\alpha_i\,
              \bigl|\tilde\lambda_i\bigr|\,\tilde K^{(i)},
  \label{eq:roeflux}
\end{equation}
%
with the wave strengths
$\tilde\alpha_{1,3} = (\Delta p \mp \tilde\rho\tilde a\,\Delta v)/2\tilde a^2$
and $\tilde\alpha_2 = \Delta\rho - \Delta p/\tilde a^2$.  The acoustic
strengths $\tilde\alpha_{1,3}$ may be formed with the velocity jump scaled by
$\min(\mathrm{Ma},1)$ after \citet{Rieper2011}, which is the published
low-Mach cure for a flux whose dissipation does not vanish with the Mach
number; the option is \texttt{Low Mach velocity jump:}, it is off by default,
and it is not recommended on a near-hydrostatic column. The factor is exactly
zero where two neighboring velocities cancel in the Roe average, so it removes
the damping of the odd-even ($2\Delta r$) velocity mode at precisely the faces
that mode lives on; measured on the LHS 1140 b low-XUV base it turned a solve
unmodified Roe completes into a nearly pure odd-even mode. Rieper's premise,
that the local and the global Mach number are comparable, does not hold on a
column running from $M\sim10^{-8}$ at the base to $10^{-2}$ at the top
(item L25, 2026-09-17).  The single index
$\bar\gamma$ is the arithmetic mean of the two face gammas, exact wherever
the two sides agree; a run whose caloric equation of state makes the index
vary across the front, that is molecular chemistry with the ladder heat
capacity, is refused this flux at startup rather than being given a wave
speed with no bound property.  A Harten-Hyman entropy correction is applied
to the acoustic eigenvalues from the star state that
\texttt{speed\_estimate\_ROE.f90} builds (the primitive-variable, two-rarefaction
and two-shock estimates of \citet{Toro2009} chapter 9, selected by the
pressure ratio).  Where that estimate has no admissible star state, two
rarefactions separating fast enough leave vacuum, or no branch returns a
positive star pressure and positive star densities, the face is given the
HLLE flux instead,
%
\begin{equation}
  F_{j+1/2} = \frac{S_R F_L - S_L F_R + S_L S_R (u_R - u_L)}{S_R - S_L},
  \qquad
  \begin{aligned}
    S_L &= \min(v_L - a_L,\ \tilde v - \tilde a),\\
    S_R &= \max(v_R + a_R,\ \tilde v + \tilde a),
  \end{aligned}
  \label{eq:hlle}
\end{equation}
%
which is positively conservative under those bounds \citep{Einfeldt1991},
where the Roe flux is not.  The number of such faces is counted and reported.

\paragraph{Local Lax-Friedrichs.}
%
\begin{equation}
  F_{j+1/2} = \tfrac12\bigl[F_L + F_R - \alpha\,(u_R - u_L)\bigr],
  \qquad
  \alpha = \max\bigl(|v_L| + a_L,\ |v_R| + a_R\bigr),
  \label{eq:llf}
\end{equation}
%
the coefficient being the spectral radius of the flux Jacobian
\citep{Rusanov1961}.  It is the most diffusive of the three and resolves no
contact wave; it is also the flux the positivity repair below substitutes at
a single face, because with cell-average input states it maps positive
density and positive internal energy to positive density and positive
internal energy under $\Delta t(|v|+c)/\Delta r \le 1$
\citep{Perthame1996, Zhang2010b}.

\subsection{Positivity repair of a stage}
\label{sec:numerics:posflux}

A Runge-Kutta stage can leave the admissible set $\rho>0$, $\rho e>0$ even
with admissible face states, because the source terms are bounded by no flux
choice.  \texttt{positivity\_limited\_fluxes} (\texttt{RK\_rhs.f90}) repairs
such a stage by replacing, at the offending cells only, the high-order
interface fluxes with the first-order Lax-Friedrichs flux of the neighboring
cell averages and rebuilding those cells' update from the mixture; the
replacement is swept until no cell is inadmissible or until both faces of a
cell are already first order, in which case no flux choice repairs it and the
caller halves the time step.  Replacing one face rather than all of them is
the flux correction of \citet{Hu2013}.  The stored face fluxes are written
back after the repair, because the transported species rows ride on
$F^\rho$ and a species flux built on a discarded high-order flux would not
sum to the mass update the density took.  Every count is reported at the end
of a run, separately for repaired faces belonging to accepted steps, so that
a run in which nothing fired is identifiable as the run an unguarded build
gives.

\subsection{Time integration and the time step}
\label{sec:numerics:rk}

The marching loop of \texttt{EXHALE\_main.f90} advances the hydrodynamics
with the three-stage strong-stability-preserving Runge-Kutta method
\citep{Gottlieb1998},
%
\begin{equation}
  \begin{aligned}
    u^{(1)} &= u^n - \Delta t\,\bigl(dF - S\bigr)\big|_{u^n},\\
    u^{(2)} &= \tfrac14\Bigl[3u^n + u^{(1)}
               - \Delta t\,\bigl(dF - S\bigr)\big|_{u^{(1)}}\Bigr],\\
    u^{n+1} &= \tfrac13\Bigl[u^n + 2u^{(2)}
               - 2\Delta t\,\bigl(dF - S\bigr)\big|_{u^{(2)}}\Bigr],
  \end{aligned}
  \label{eq:ssprk3}
\end{equation}
%
each stage followed by the boundary condition, by the admissibility test of
Sec.~\ref{sec:numerics:posflux}, and by the species advection sub-step that
rides on the face mass fluxes that stage built.  A stage that cannot be
repaired causes the whole step to be discarded and retaken from $u^n$ at half
$\Delta t$, up to a bounded number of bisections.

The time step is the explicit-stable interval of the scheme actually in use,
derived at \texttt{eval\_dt.f90} for the spherical finite-volume divergence
the right-hand side assembles,
%
\begin{equation}
  dF_j = \frac{A_{j+1/2}F_{j+1/2} - A_{j-1/2}F_{j-1/2}}{V_j},
  \qquad A = r^2, \qquad
  V_j = \frac{r_{j+1/2}^3 - r_{j-1/2}^3}{3}.
  \label{eq:fvdiv}
\end{equation}
%
Every stage of \eqref{eq:ssprk3} is a convex combination of forward-Euler
steps of the same length (the SSP coefficient is 1), so the restriction is
the forward-Euler one of \eqref{eq:fvdiv}.  Each face flux is a Riemann flux
of the two states bounding that face, and the fastest signal of that face's
Riemann problem is bounded by
$S_{j+1/2} = \max(|v|+c_s)$ over the two cell averages bounding it.  Godunov's
condition is that the fan opened at a face has not left the cell it enters
when the step ends; written with the geometry \eqref{eq:fvdiv} uses, the face
carries signal into cell $j$ at the rate $A_fS_f$ per unit volume $V_j$, so
%
\begin{equation}
  \Delta t = C_{\rm CFL}\,\min_{1\le j\le N}
  \frac{V_j}{\max\bigl(A_{j-1/2}S_{j-1/2},\ A_{j+1/2}S_{j+1/2}\bigr)},
  \qquad C_{\rm CFL} = 0.6 \ \text{by default},
  \label{eq:cfl}
\end{equation}
%
with $c_s = \sqrt{\gamma_{\rm eff}p/\rho}$ from the cell's own composition.
In the Cartesian uniform limit $A_{j-1/2}=A_{j+1/2}$ and $V_j = A\,\Delta r$,
and \eqref{eq:cfl} returns $C_{\rm CFL}\Delta r/(|v|+c_s)$, so the number the
\texttt{CFL:} key carries keeps its usual meaning and nothing is
recalibrated.  The stricter condition that the two fans of one cell may not
meet inside it replaces the maximum by the sum of the two face terms, which
is the same restriction at $C_{\rm CFL}/2$.

Two properties of \eqref{eq:cfl} are worth stating, because the cell-centered
form $\Delta r_j/(|v_j|+c_{s,j})$ has neither.  In spherical geometry
$V_j/A_{j+1/2} < \Delta r_j < V_j/A_{j-1/2}$, so the cell-centered width is
not an upper bound of the interval on either face; on the production base
grid the two differ by $\Delta r/r = 1.96\times10^{-4}$ (MEASURED, item P6b).
And the minimum is taken over the evolved cells $1\le j\le N$ alone.  A ghost
row is overwritten by the boundary condition after every stage, so its own
width is not a stability requirement, and the lowest ghost does not even
carry a width of its own; the ghosts nevertheless bound the two boundary
faces $r_{\rm edg}(0)$ and $r_{\rm edg}(N)$, which are faces of cells 1 and
$N$, so the base Riemann problem restricts cell~1 with the reservoir's own
signal speed although no ghost cell is evolved.  Letting the ghosts carry
widths of their own is not conservative in the other direction either: an
outer ghost given a sound speed a thousand times the column's shortened the
step of a 500-cell domain by 89.7 per cent (MEASURED, item P6b).

$S_f$ is formed from the cell averages bounding the face and not from the
reconstructed face states, which inside a steep cell can carry a larger
$|v|+c_s$; that margin is the one a Godunov number below unity leaves.
With \texttt{Time stepping: Local} each cell instead takes its own
$\Delta t_j$, the same expression without the minimum; since the fixed point
$dF = S$, $\mathcal{H} = \Lambda$ is independent of the step, this is a
legitimate acceleration toward a steady state, and it makes the trajectory
meaningless as a time history.  In the global mode the local array is
uniformly the global step, which keeps the update arithmetic identical to the
scalar form.  A ghost row, which is not evolved, is given the interval of the
evolved cell it borders so that a stage formed over the padded range is
arithmetic on a physical interval.

\eqref{eq:cfl} is not only a marching quantity: it supplies the stationary
route's initial pseudo-time and the floor of the JFNK pseudo-transient
continuation (Sec.~\ref{sec:steady:ptc}), so its value moves solver behavior
as well as a time step.  The no-step evaluation route does not call it.

\subsection{Operator splitting of one step}
\label{sec:numerics:split}

One accepted marching step applies the following operators, in this order,
with the state after each one carried into the next
(\texttt{EXHALE\_main.f90}):
%
\begin{enumerate}
\item the three Runge-Kutta stages \eqref{eq:ssprk3}, each with its
      reconstruction, flux assembly, boundary fill, positivity test and
      species advection sub-step;
\item the conversion to primitive variables, the composition-dependent
      particle and electron counts, and the temperature
      \eqref{eq:thermal_eos};
\item the advected element composition written into the species vector
      (\texttt{species\_advection\_project}), followed by the implicit
      diffusive half of the element transport
      (\texttt{element\_diffusion\_step} in
      \texttt{binary\_element\_diffusion.f90}, Sec.~\ref{sec:mol:diffusion}),
      active only with \texttt{He\_diffusion}, and, on an accepted element
      step, the energy update $\mathcal{E} \leftarrow \mathcal{E} - \Delta
      t\,\nabla\!\cdot q_d$ from the element face fluxes of that step
      (Sec.~\ref{sec:mol:interdiffusion});
\item the vertical transport of the carriers with their chemistry
      (\texttt{photochemical\_transport\_step} in
      \texttt{diffusive\_photochemistry.f90}, Sec.~\ref{sec:mol:carrier}):
      the molecular carriers with \texttt{Molecular carrier transport}, the
      ionization stages with \texttt{Ionization transport} and He\,$2^3$S with
      \texttt{He 2\^{}3S transport}, followed, where that interval was
      covered, by $\mathcal{E} \leftarrow \mathcal{E} - \Delta
      t\,\nabla\!\cdot q_c$ from the composition the step returned
      (\texttt{carrier\_enthalpy\_divergence\_of\_state});
\item the refresh of the lagged $\mathrm{H}(n{=}2)$ populations and their
      Balmer source (\texttt{excited\_hydrogen.f90},
      Sec.~\ref{sec:ionization:n2});
\item the coupled source step: a bounded fixed-point iteration alternating
      the composition equilibrium at the current temperature
      (\texttt{ioniz\_eq}, Sec.~\ref{sec:ionization:solve}) with the
      semi-implicit energy update \eqref{eq:th:semiimplicit} at the returned
      composition, exiting when the temperature and the composition both stop
      moving;
\item the boundary condition;
\item the Crank-Nicolson viscous and conductive stage
      \eqref{eq:viscous}, where those options are on, and the boundary
      condition again;
\item the periodic Shapiro filter, where it is on, and the boundary
      condition again.
\end{enumerate}
%
Steps 3, 4, 8 and 9 are no-ops unless their key is present, and the two
enthalpy updates also need \texttt{Interdiffusion enthalpy flux} (on by
default).  In a physical-mode run (a trajectory with a clock) a refused element
step or a carrier interval the transport operator could not cover refuses the
whole attempted step, which is restored from its checkpoint and retaken
with a reduced step (\texttt{attempted\_step\_reduced\_dt}) up to a retry
budget whose exhaustion stops the run; while
the run is reaching a state the refusal is counted and the step goes on.  A
step whose energy update found no temperature, or whose Crank-Nicolson stage
failed (a temperature below the floor of the equation of state, a non-finite
solution, a pivot breakdown; \texttt{viscous\_conduction\_step} prints the
cell, the solved value and the energy a floor would have injected), has no
state to march on and is refused in both modes: it is restored and retaken
at half the interval within the same budget, whose exhaustion stops the run
with exit status~2 (Sec.~\ref{sec:hydro:transport}).  The Shapiro
filter is the 1-2-1 low-pass pass
$u_j \leftarrow u_j + \tfrac14\epsilon(u_{j-1} - 2u_j + u_{j+1})$ applied
every \texttt{shapiro\_every} steps, the way CETIMB damps the base sound wave
\citep{Koskinen2013}.  It is off by default and is recorded as a validity
state of every step it runs on, because unlike the low-Mach term of
Sec.~\ref{sec:hydro:lowmach} it alters the marched state with no source term
behind it and therefore has no counterpart in the stationary residual.

\subsection{The marching convergence measure and the stop conditions}
\label{sec:numerics:du}

The measure the marching stops on is the radial spread of the mass flux over
the escape window,
%
\begin{equation}
  \mathrm{du} = \frac{\max_{j\ge j_{\min}} F_j - \min_{j\ge j_{\min}} F_j}
                     {\bigl|\,\overline{F}\,\bigr|},
  \qquad F_j = \rho_j v_j r_j^2,
  \label{eq:du}
\end{equation}
%
with $\overline F$ the mean over the same window.  A steady spherical wind
carries one mass flux at every radius, which is the stop ATES uses
($\Delta\dot M/\dot M < 10^{-3}$) and the constancy of $F_c = \rho v r^2$
CETIMB asks for.  The sign is kept and the normalization is the mean: taking
$|F|$ first would read a window whose flux reverses, an inflow at one radius
and an outflow at another, as flat.  It is evaluated once per step and read
by every consumer, so the stop, the trigger arming, the staged
secondary-ionization activation of Sec.~\ref{sec:radiation:secondary} and the
handoff to the steady solver all speak about one number.  This is emphatically \emph{not} a velocity change, and a single-stage PLM run
stalls at a few percent spread.

A second measure, the largest relative change of the conserved state over a
step,
%
\begin{equation}
  \mathrm{dtu} = \max_{k\in\{1,2,3\}}\ \max_{j_{\min}\le j\le N}
                 \left|1 - \frac{u^{n+1}_{k,j}}{u^{n}_{k,j}}\right|,
  \label{eq:dtu}
\end{equation}
%
says that the state has stopped changing; a row that was exactly zero at the
start of the step has no relative change defined and the measure is taken as
unbounded rather than dropped.

A run stops when one of the following holds, and reports which:
$\mathrm{du} < \texttt{du\_th}$ (default $10^{-3}$) with the trigger armed
and the flux level stable; $\mathrm{dtu} < \texttt{dtu\_th} = 10^{-8}$;
\texttt{du} settled on a plateau over \texttt{N\_stall} steps (default 2000)
within \texttt{stall\_tol} (default $10^{-6}$); the steady gates of
Sec.~\ref{sec:steady:gates} where \texttt{Resid tol:} was stated; the Newton
finish returning a state; the step cap \texttt{EXHALE\_MAXSTEPS} where it is
set (a relaxation snapshot, which the regression cases that are not converged
solutions pin in their \texttt{maxsteps} file); or the hard cap
\texttt{count\_max} ($10^6$).  The
triggers are \emph{armed} the first time the measure is seen at or above its
threshold, so only a descending crossing can fire them and a run whose initial
condition already carries a flat mass flux, which a transonic one does, cannot
stop at step zero.  With \texttt{Solver: Newton} the marching is a warm-up
only: the du and dtu stops are disabled and the loop hands over once
\texttt{du} falls below \texttt{newton\_du\_switch} (default $10^{-2}$) or
plateaus near it.  Both measures can be small at a premature balance of the
operator split while the energy residual is of order 30, so for a quantitative
mass-loss rate the stationary finish is not optional.  Every thousandth step
the marching loop writes the state it holds; that write is a snapshot and not
the run's answer, and carries \texttt{certified=F
cert\_reason=marching\_snapshot\_s<step>}, the run's own certificate being
restored after it, because the last certificate the run measured belongs to
another state.
